% !TEX root = ../main.tex
% ======================================================================
% Section 5 – Analyse d'intégrité, éthique et ÉDI
% Grille : « Analyse éthique & ÉDI » – 16 pts (section prioritaire).
% Longueur indicative : ~4 à 5 pages.
% Cadres obligatoires : \eptc{} (dont le chapitre 9) et la \loivingtcinq{}.
% ======================================================================
% FIL CONDUCTEUR (toute la section) : analyse critique de la promesse
% « les données ne quittent jamais l'établissement ».

\subsection{Fil conducteur : la promesse « les données ne quittent jamais l'établissement »}
\label{sec:promesse}
% TODO : ce que la promesse garantit réellement et ce qu'elle ne couvre pas.
\arediger{analyse critique de la promesse}

\subsection{Intégrité et éthique}
\label{sec:integrite-ethique}

\subsubsection{Risques résiduels}
\label{sec:risques-residuels}
% TODO : quels risques résiduels subsistent, dont l'inférence d'appartenance et l'inversion de modèle à partir des mises à jour ?
% TODO : la confidentialité différentielle couvre-t-elle ces risques ? À quel prix (compromis ε / performance) ?
\arediger{risques résiduels}

\subsubsection{Choix du modèle de consentement}
\label{sec:consentement}
% TODO : quel modèle de consentement choisir : large, spécifique ou dynamique ? Justifier au regard de l'EPTC 2 et de la Loi 25.
\arediger{consentement}

\subsubsection{Usages secondaires du modèle}
\label{sec:usages-secondaires}
% TODO : quels usages secondaires du modèle sont envisageables, et lesquels doivent être encadrés ou interdits ?
\arediger{usages secondaires}

\subsubsection{Proposition d'un protocole de gouvernance}
\label{sec:protocole-gouvernance}
% TODO : proposer un protocole de gouvernance (instances, rôles, décisions, rôle des CER et des instances autochtones).
\arediger{protocole de gouvernance}

\subsection{Équité, diversité et inclusion (ÉDI)}
\label{sec:edi}

\subsubsection{Agrégation pondérée et petits sites}
\label{sec:edi-petits-sites}
% TODO : l'agrégation pondérée par la taille marginalise-t-elle les petits sites ?
\arediger{petits sites}

\subsubsection{Performance par sous-groupe}
\label{sec:edi-sous-groupes}
% TODO : le modèle performe-t-il également pour chaque sous-groupe ? Comment le mesurer ?
\arediger{sous-groupes}

\subsubsection{Prévention de la stigmatisation}
\label{sec:edi-stigmatisation}
% TODO : comment éviter la stigmatisation des communautés plus touchées par le diabète ?
\arediger{stigmatisation}
